%=============================================================================!
% 13.0  CHAPTER OPENING                                                       !
%=============================================================================!
Chapter~12 described the canonical distribution of a system in contact
with a heat bath. To reproduce that situation in a simulation, an
algorithm must exchange energy with the atoms while preserving the
required distribution. Such an algorithm is called a thermostat. We
judge it by two questions: whether it samples the canonical ensemble,
and how much it changes the motion during sampling.

We develop six approaches, beginning with direct velocity rescaling and
ending with the stochastic rescaling used by TrajCast. Comparisons on
liquid argon and on lithium atoms moving over the model surface of
Chapter~3 show why a suitable method depends on the measurement. The
chapter draws together choices for equilibration, equilibrium averages
and dynamics, then examines the signs of faulty thermostatting in a
trajectory. Chapter~14 applies the same distinction between mean values,
fluctuations and motion to pressure control.
